\section{A uniform obstruction for scalar exterior-square permutation blocks}
\label{sec:exterior-square-obstruction}

\paragraph{Scope and check status.}
This is a new elementary proof, initially recorded on 21 November 2026.
It rules out a dominant sixfold source before solving the product-one
equations: even its necessary branch profiles do not exist.  The exact
partition calculation accompanying the proof checks the spectral formulas
for every partition through degree forty at scalar orders three, four,
and six (conjugate phases have equal ranks), and checks them against direct
exterior-square eigenvalues through degree twelve.  That finite calculation
is a check of the formulas, not a substitute for the uniform proof below.
The parent and the independent G34 agent have each completed a whole-proof
review, including the projection-descent and automatic-transitivity
arguments, and recorded PASS.  The parent reran the full finite certificate;
the G34 agent independently rebuilt the fractional spectra and signed pair
orbits through degree twelve without importing this certificate's code.
The only general geometric inputs are the standard finite-cover Euler
formula, Riemann--Hurwitz, and the previously proved multiplicity-variation
factorization in D-BI-01.  No classification theorem is used here.

\subsection*{D-EX-01. Three uniform spectral inequalities}

Let $n\ge4$, let $U$ be the deleted permutation representation of
$\mathfrak S_n$, and put
\[
        W=\bigwedge^2 U,\qquad m=\dim W=\binom{n-1}{2}.
                                                               \tag{EX.1}
\]
For a permutation $\pi$, write $c$ for its number of cycles, including
fixed points, and $k=n-c=\operatorname{ind}(\pi)$.  For a complex root
of unity $\lambda$, put
$D_\lambda(\pi)=\operatorname{rank}(1-\lambda W(\pi))$.
We prove
\begin{align}
 \lambda\ne1&\quad\Longrightarrow\quad
 D_\lambda(\pi)\ge \frac{n-4}{2}k+1,                  \tag{EX.2}\\
 \lambda\ne1&\quad\Longrightarrow\quad
 D_\lambda(\pi)\ge
 A_n:=m-\left\lfloor\frac{(n-1)^2}{4}\right\rfloor,  \tag{EX.3}\\
 \pi\ne1&\quad\Longrightarrow\quad
 D_1(\pi)\ge\frac{n-4}{2}k+\frac n3.                 \tag{EX.4}
\end{align}
These estimates hold for all scalar orders, not only orders three, four,
and six.

\paragraph{Proof of (EX.2).}
Put $t=\lambda^{-1}\ne1$.  Every eigenvalue of the permutation
representation has multiplicity at most $c$.  On $U$, the eigenvalue
$1$ has multiplicity $a=c-1$, and every other eigenvalue still has
multiplicity at most $c$.  Its eigenvalue-$1$ complement has dimension
$k$.  Let $b$ be the multiplicity of $t$ on that complement.
The fixed dimension in question is the multiplicity of $t$ on
$\bigwedge^2U$.

First suppose $2\le c\le n/2$.  Pair eigenvalue labels by the
involution $z\mapsto t/z$.  A two-element pair with multiplicities
$x,y\le c$ contributes $xy\le c(x+y)/2$; a fixed label contributes
$\binom{x}{2}\le cx/2$.  Treat the pair $1,t$ separately: it contributes
$ab$.  Thus the desired multiplicity $F$ satisfies
\[
 F\le ab+\frac c2(n-1-a-b)
   \le\frac c2(n-2).
\]
For the last inequality, use $a=c-1$, $b\le c$, and
\[
 \frac c2(a+b)-ab
       =\frac{c(c-1)-(c-2)b}{2}\ge\frac c2.
\]
Consequently
\[
 D_\lambda(\pi)-\left(\frac{n-4}{2}k+1\right)
       \ge k-\frac n2\ge0.
\]
If $c=1$, all eigenvalues on $U$ are distinct.  Pairing their labels
shows $F\le\lfloor(n-1)/2\rfloor\le n-2$, which is precisely enough
for (EX.2), since $k=n-1$.

It remains to treat $c>n/2$.  Here $a=c-1\ge k$.  In the complement
to the eigenvalue-$1$ space, an eigenvector with eigenvalue $t$ cannot
participate in an exterior pair of product $t$: its partner would
have eigenvalue $1$.  It follows that
\[
 F\le ab+\binom{k-b}{2}\le ak.
\]
Indeed the middle expression is a convex quadratic in $b\in[0,k]$,
and its endpoint values are $\binom k2$ and $ak$, the latter being
larger.  Since $n\ge4$ gives $c\ge3$ in this case,
\[
 D_\lambda(\pi)-\left(\frac{n-4}{2}k+1\right)
  \ge \frac{(c-k)(c-3)}2\ge0.
\]
This proves (EX.2).

\paragraph{Proof of (EX.3).}
The representation $W$ is real.  If $\lambda$ is nonreal, the
eigenvalues $\lambda^{-1}$ and $\overline\lambda^{-1}$ on $W$ have
equal multiplicity, so each has multiplicity at most $m/2$, and hence
at most $\lfloor(n-1)^2/4\rfloor$.
The only remaining nonidentity real scalar is $\lambda=-1$.
On $U$ pair eigenvalue labels by $z\mapsto-1/z$.
Each nonfixed pair of multiplicities $x,y$ contributes $xy$ to
the eigenvalue-$(-1)$ space in $\bigwedge^2U$, at most
$(x+y)^2/4$.  The two fixed labels are $i,-i$; their multiplicities
are equal, say $b$, because $U$ is real.  Taken together they
contribute $2\binom b2=b(b-1)\le(2b)^2/4$.
These blocks partition a space of total dimension $n-1$.
Summing their square bounds gives
\[
 \dim\ker(1+W(\pi))
       \le\frac14\sum_{\text{blocks}}(\dim\text{block})^2
       \le\frac{(n-1)^2}{4}.
\]
Taking the integer part proves (EX.3).

\paragraph{Proof of (EX.4).}
Let $l_1,\ldots,l_c$ be the cycle lengths, and let $e$ count the
even lengths.  On the exterior square of the full permutation space,
two different cycles contribute $\gcd(l_i,l_j)$ fixed vectors.
Inside one cycle of length $l$, the contribution is
$\lfloor(l-1)/2\rfloor$: the opposite-pair orbit for even $l$ has
negative return sign and contributes no fixed vector.  Since
$\bigwedge^2(1\oplus U)=U\oplus\bigwedge^2U$, subtraction of the
$c-1$ fixed vectors of $U$ gives the exact identity
\[
 D_1(\pi)=\frac{(n-3)k+E}{2},\qquad
 E=e+\sum_{i<j}\bigl(l_i+l_j-2\gcd(l_i,l_j)\bigr).
                                                               \tag{EX.5}
\]
All summands of $E$ are nonnegative.  We claim $k+E\ge2n/3$ for
every nonidentity permutation.
If there are no fixed points, let $o$ count the odd cycles.
They all have length at least three, so $o\le n/3$, and
$k+E\ge n-c+e=n-o\ge2n/3$.
If there are $f\ge1$ fixed points, the cross terms involving them
give $E\ge fk$.  Also $k\ge1$ and $n\le f+2k$, since each
nontrivial cycle satisfies $l\le2(l-1)$.  Hence
\[
 3(k+E)\ge3(f+1)k\ge2f+4k\ge2n;
\]
the middle inequality is $(3f-1)k\ge2f$, valid for $f,k\ge1$.
Combining this claim with (EX.5) proves (EX.4).

\subsection*{D-EX-02. No twelve-branch signature-$(3,3)$ profile}

Suppose $n\ge5$ and a product-one tuple of genuine branch elements
$g_j=\lambda_j W(\pi_j)$ has transitive permutation projection.
Assume its sixfold multiplicity satisfies the necessary genus-zero
Euler equation
\[
             \sum_jD_{\lambda_j}(\pi_j)=2m+6.            \tag{EX.6}
\]
For completeness, the general Euler formula follows by extending the
rank-$m$ local system over each puncture with its inertia invariants.
Its Euler characteristic on a genus-$h$ curve is
$m(2-2h-r)+\sum_j\dim\ker(g_j-1)$.
For a nontrivial irreducible local system its degree-zero and degree-two
cohomology vanish, so its first cohomology has dimension
$2m(h-1)+\sum_jD_j$.  Setting this dimension equal to six and $h=0$
gives (EX.6).
The transitive permutation cover of degree $n$ has nonnegative genus,
so Riemann--Hurwitz gives
\[
                    \sum_j\operatorname{ind}(\pi_j)\ge2n-2.
                                                               \tag{EX.7}
\]
We prove that the tuple has at most eleven entries.

Suppose instead that $r\ge12$, and let $s$ count the entries
with $\lambda_j\ne1$.  For $n\ge6$, (EX.3) and (EX.6) imply
$s\le6$.  Here is the full numerical check, valid for every $n$.
Writing $n-1=2v+1$, with $v\ge2$, gives $A_n=v^2$ and
\[
             7A_n-(2m+6)=3v^2-2v-6>0.
\]
Writing $n-1=2v$, with $v\ge3$, gives $A_n=v(v-1)$ and
\[
             7A_n-(2m+6)=3v^2-5v-6>0.
\]
Every scalar-$1$ entry has $\pi_j\ne1$, because branch entries are
genuine.  Thus (EX.2), (EX.4), and (EX.7) yield
\begin{align*}
 \sum_jD_{\lambda_j}(\pi_j)
 &\ge\frac{n-4}{2}\sum_j\operatorname{ind}(\pi_j)
            +\frac n3(r-s)+s\\
 &\ge(n-4)(n-1)+\frac n3(12-s)+s\\
 &\ge(n-4)(n-1)+2n+6
       =n^2-3n+10=2m+8.
\end{align*}
The last inequality uses $n\ge6$ and $s\le6$.
This contradicts (EX.6).

For $n=5$, one has $m=6$, $A_n=2$, and $s\le9$.
The quantity $D_1(\pi)-k/2$ belongs to $\frac12\mathbf Z$;
(EX.4) bounds it below by $5/3$, and hence by $2$.
Using (EX.2) for the other entries gives
\[
 \sum_jD_{\lambda_j}(\pi_j)
       \ge\frac12(2n-2)+2(r-s)+s
       \ge28-s\ge19>18=2m+6.
\]
This handles $n=5$ as well.

For $n=4$, the same conclusion holds whenever the scalar quotient
is nonreal.  In this case $m=3$ and (EX.6) is $\sum_jD_j=12$.
There must be at least two nonreal scalar entries: there is at least
one because the entries generate the nonreal scalar image, and the
product of one nonreal scalar with any collection of $\pm1$ cannot
be one.  Since $W$ is real of dimension three, each of these entries
has moving dimension at least two.  Every other genuine entry has
moving dimension at least one.  Therefore
$\sum_jD_j\ge r+2\ge14$, again a contradiction.

\subsection*{D-EX-03. Consequence for inducing blocks and the exact spectrum}

Let an exact-imaginary-quadratic-field sixfold representation be
induced from an absolutely irreducible representation $\tau$ of
dimension $m=\binom{n-1}{2}$, whose finite image in this representation
has the form
\[
      \tau(j)=\chi(j)\,\bigwedge^2\operatorname{Std}_n(\pi(j)),
      \qquad n\ge4,
                                                               \tag{EX.8}
\]
There is no restriction on the order or coefficient field of $\chi$;
the permutation image need not be all of $\mathfrak S_n$.
Then its period rank is at most eight; in particular it supplies
no dominant Weil sixfold family, of either determinant class.

The permutation image is automatically transitive.  Otherwise the
deleted permutation space contains a nonzero invariant vector.
By averaging a Hermitian form, split off its invariant line as
$U=1\oplus U'$, where $\dim U'=n-2\ge2$.  Then
$\bigwedge^2U=U'\oplus\bigwedge^2U'$ is a direct sum of two
nonzero invariant subspaces.  A scalar twist preserves these
subspaces, contrary to irreducibility of $\tau$.

Both the scalar and permutation maps descend when the kernel of
$\tau$ is quotiented.  Indeed $\bigwedge^2\operatorname{Std}_n$
is projectively faithful for $n\ge4$: if $\bigwedge^2T$ is scalar
for a semisimple $T$ on a space of dimension at least three, the
equalities between all pairwise eigenvalue products force every
eigenvalue of $T$ to be the same.  Thus $T$ itself is scalar.
The deleted permutation representation contains no nonidentity
scalar.  One way to see the latter assertion is to apply the
permutation to every difference $e_i-e_j$: scalar action preserves
every two-element support $\{i,j\}$, which forces the permutation
to be the identity for $n\ge3$.  Therefore an element of
$\ker\tau$ has $\pi=1$ and then $\chi=1$.

To justify passage to the intermediate cover, quotient the kernel
of $\tau$ in its inducing subgroup.  Frobenius reciprocity identifies
the complex six-dimensional multiplicity variation with the
$\tau$-multiplicity on the resulting finite cover of a curve $X$.
As in D-BI-01, the fixed Hermitian polarization determines the
conjugate Hodge subspace from the chosen one.  Thus the full
quadratic-field period rank factors through this intermediate-cover
moduli even when $\chi$ takes values in a larger cyclotomic field.
All finite-cover and inducing data over a fixed pointed $X$ are
discrete.

If $h=g(X)$, the Euler equation is
\[
                        6=2m(h-1)+\sum_jD_j.
                                                               \tag{EX.9}
\]
For $h=1$ there are at most six genuine marked points, and the
marked-base moduli dimension is at most six.  For $h\ge2$, since
$m\ge3$, the only possible case is $h=2,m=3$ with no ramification;
its base moduli has dimension three.  Rank nine would therefore
require $h=0$ and at least twelve branch points.  This is excluded
by D-EX-02.  In the $n=4$ case the scalar image is necessarily
nonreal: if it were contained in $\{1,-1\}$, $\tau$ would have
real character and so would its induced representation, contrary
to its exact imaginary quadratic character field.  This proves
the period-rank assertion.

For clarity, the complete finite partition spectrum used for checking
these estimates is also explicit.  For a partition
$l_1+\cdots+l_c=n$, the standard-representation eigenvalues are the
fractions
\[
 \coprod_i\{0,1/l_i,\ldots,(l_i-1)/l_i\}
                  \quad\text{with one zero removed}.
\]
If $\lambda=\exp(2\pi i a/d)$, add $a/d$ to every unordered sum
of two distinct members of this multiset, and reduce modulo one.
The number of nonzero resulting fractions is $D_\lambda$; their
sum is the age.  Signature $(3,3)$ additionally requires total
age $m+3$, as well as (EX.6).  Since (EX.6)--(EX.7) and $r\ge12$
already contradict the inequalities, imposing these extra
Chevalley--Weil, scalar-product, and tuple-generation constraints
cannot restore a profile.

An equivalent integer formula for the fixed dimension is useful
for the certificate.  Let $e$ be the order of $\lambda$.  On
$\bigwedge^2$ of the full permutation space, a pair of different
cycles contributes
\[
       \gcd(l_i,l_j)\,\mathbf1_{e\mid\operatorname{lcm}(l_i,l_j)}.
\]
A single cycle of length $l$ contributes
\[
  \left\lfloor\frac{l-1}{2}\right\rfloor\mathbf1_{e\mid l}
       +\mathbf1_{\,2\mid l\,,\ \lambda^{l/2}=-1}.
\]
Subtract $\sum_i\mathbf1_{e\mid l_i}$, and add $1$ if
$\lambda=1$, to pass to $\bigwedge^2U$.
This formula counts signed pair orbits, while the direct fraction
calculation counts eigenvectors; the executable check compares the
two without floating-point arithmetic.
